% ---------- ABSTRAK ----------
% Jumlah halaman "x+N" dihitung otomatis: x = halaman romawi terakhir,
% N = halaman isi terakhir (lihat \label{akhirfront} & paket lastpage).
\chapter*{Abstrak}
\addcontentsline{toc}{chapter}{ABSTRAK}
\begin{spacing}{1.5}
\noindent\textbf{\skNama}, ``\skJudul''.\par
\noindent\pageref{akhirfront}+\pageref{LastPage} halaman\par
\vspace{0.5\baselineskip}

\noindent\justifying
Abstrak berisi ringkasan skripsi yang dideskripsikan secara jelas, padat, dan
singkat, dengan maksimum 200 kata. Isi abstrak terdiri atas permasalahan
penelitian, deskripsi metode/sistem yang dikembangkan, hasil evaluasi, kesimpulan,
dan saran. Abstrak ditulis dalam satu atau dua paragraf, diketik 1,5 spasi, serta
menghindari sitasi dan singkatan tak baku.
\vspace{0.5\baselineskip}

\noindent\textbf{Kata Kunci:} \hangindent=2.2cm kata kunci pertama, kata kunci kedua,
kata kunci ketiga, hingga maksimum enam kata kunci.
\end{spacing}
